
\documentclass[submit,autodetect-engine,dvipdfmx,platex]{ipsj-mod}

\usepackage{amsmath}
\usepackage{graphicx}       % 画像を挿入できるようにする
\usepackage{wrapfig}        % 図の回り込みを許す
\usepackage{lineno}         % 行番号をつける
\usepackage{xcolor}         % 文字に色をつけられるようにする
\usepackage{colortbl}       % 表に色をつけられるようにする
\usepackage{booktabs}       % 表に横罫線をつけられるようにする
\usepackage{multirow}       % 複雑な表を書けるようにする
\usepackage{ascmac}         % 枠をつけた文を書けるようにする
\usepackage{latexsym}       % 扱える数学記号を追加
\usepackage{bm}             % ベクトル表記を扱えるようにする
\usepackage{array}          % 数式モードで表（や行列）を書けるようにする
\usepackage{hyperref}       % URLをハイパーリンク化する
\usepackage{pxjahyper}      % URLのハイパーリンク化した際の日本語問題を解決する
\usepackage{ccicons}        % クリエイティブ・コモンズ・ライセンスのマークをつけられるようにする
\usepackage[utf8]{inputenc} % for a UTF8 editor only
%\usepackage{microtype}      % Improved Tracking and Kerning

% hyperrefが\cite番号をリンク化するとipsjクラスの引用番号ソート処理と衝突するため、
% \bibciteを標準の定義（番号のみを保存）に戻す。URLや\hrefのリンクは従来どおり有効。
\makeatletter
\def\bibcite#1#2{\@newl@bel{b}{#1}{#2}}
\makeatother

\def\Underline{\setbox0\hbox\bgroup\let\\\endUnderline}
\def\endUnderline{\vphantom{y}\egroup\smash{\underline{\box0}}\\}
\def\|{\verb|}

\begin{document}


% ----------------------------
% 論文のタイトル
% ----------------------------
\title{論文のタイトル}

% ----------------------------
% 著者の情報
% ----------------------------
\author{山畑 滝子}{Takiko Yamabata}{NCU}[yamabata@ed.nagoya-cu.ac.jp]
\affiliate{NCU}{名古屋市立大学 データサイエンス学部}



% ----------------------------
% アブストラクト（300〜500字程度）
% ----------------------------
\begin{abstract}
% -----------------------------
% 執筆にあたっての留意点
% -----------------------------
% 実際にアブストラクトを執筆する際には，以下の項目を順番に1〜2行ずつ書くと良いでしょう．

% 1. 論文の目的（何を提案する，何を調査するのか）の要約
% 2. 提案手法・調査事項の具体的なアイデア
% 3. 評価実験・調査の結果，明らかになったこと
% 4. 本論文によってどんなことが可能になるか，社会や学術コミュニティにどんな波及効果があるか
%
%
% -----------------------------
% 補足（参考：how-to-xx/how-to-write-a-paper/1-学術論文の書き方.md 3.2節）
% -----------------------------
%
% ■ アブストラクトは論文全体のミニチュアです
% 読者や査読者の多くは，「アブストラクト → 図 → 結論 →（気になれば）本文」という順に流し読みします．
% つまり，多くの人にとってアブストラクトが唯一読まれる部分です．
% ここだけを読んでも，何を問題とし，何を提案し，何がわかったのかが伝わる状態を目指してください．
%
% ■ 背景や動機を長々と書いてはいけません
% 学生が最もやりがちな失敗は，背景の説明にアブストラクトの半分を使ってしまうことです．
% 背景に割いてよいのは，多くとも1文です．上の項目1〜4のうち，紙面の大半は項目2と項目3に使ってください．
% 分量の目安は次の通りです．
%
% * 背景：0〜1文
% * 問題：1文
% * 既存手法の限界：1文
% * 提案：1〜2文
% * 結果：1〜2文
%
% 和文アブストラクトは300〜500字が目安です．
% この分量の目安に沿って1〜2文ずつ書いていけば，おおむねこの範囲に収まります
% （下の例は約440字です）．
%
% ■ 結果は，可能な限り数値で書きます
% 「大幅に改善した」ではなく「従来手法と比べて◯%改善した」「検索回数が◯倍に増えた」と書くと，
% 主張が一気に強くなります．数値のないアブストラクトは，読者に「本当に効果があったのか」を判断させられません．
%
% ■ 本文の文を一語一句そのままコピーして使わないでください
% アブストラクトはイントロダクションの要約ではなく，独立した一つの文章として書きます．
% イントロダクションの第1段落をそのまま貼り付けたものは，すぐに見抜かれます．
%
% ■ その他の作法
% * 原則として，アブストラクトには文献を引用しません（引用がないと意味の通らない書き方を避けます）
% * 略語は，アブストラクト内で初めて使うときに展開します
%   （本文でも改めて展開します．アブストラクトと本文は，初出を別々に数えます）
% * 図表や数式への参照（「図1に示すように」など）も書きません
% * 書くのは最後です．本文と結論を書き終えてから，論文全体を眺めて書くのが最も速く，正確です
%   （ただし「研究を始めた時点から書き始める論文」では，現時点で書けることを書いて，何度も更新していきます）
%
%
% 以下は例です．
% 各文末の「（項目◯に対応）」は，上の項目との対応を示すための注記です．
% 自分の論文を書く際には削除してください．

本稿では，ウェブ検索中のユーザに注意深い情報探索を暗黙的に促す「クエリプライミング」を提案する（項目1に対応）．
クエリプライミングは，批判的思考を喚起し注意深い情報探索や意思決定を促進するようなキーワードを，検索ワード補完・推薦時に提示する（項目2に対応）．
クエリプライミングの有効性を検証するために，クラウドソーシングを用いたオンラインユーザ実験を行った．
被験者の情報探索ログ分析および実験のアンケート調査の結果，以下の傾向が明らかになった：
(1) クエリプライミングが実装されたウェブ検索エンジンを用いた被験者はセッション中の検索回数が増え，検索結果一覧ページを何度も見直すようになる．
(2) クエリプライミングによって，証拠を重視してウェブページを検索・閲覧する行動が促進される．
(3) クエリプライミングの効果は被験者の学歴に依存する（項目3に対応）．
本研究で得られた知見によって，注意深い情報探索を活性化させる検索インタラクションの設計に寄与することが期待される（項目4に対応）．

\end{abstract}

% ----------------------------
% 研究のキーワード
% 5つ程度のキーワードをカンマ（，）で区切る
% ----------------------------
\begin{jkeyword}
情報検索，意思決定支援，生成AI
\end{jkeyword}


\maketitle


% ------------------------------
% メインコンテンツ
% ------------------------------

% 1章：はじめに
% -----------------------------
% 執筆にあたっての留意点
% -----------------------------
%
% 論文の第1章にあたる文章です．イントロダクションは極めて重要です．イントロダクションの役割は，読者に論文を読むべきか否かを敏速・的確に判断するための材料を提示すること，論文を読むために必要となる前提知識を提供することにあります．読者はイントロダクションを読んで面白いと思えば論文を読み続けるし，面白くなければ論文を読むのをやめます．イントロダクションで読者を惹きつけられなければ，そこで試合終了です．
%
% イントロダクションでは，以下の要素を順番に書いていくと良いでしょう．
%
% #### 手法・システム提案系論文の場合
% 1. （次に述べる問題につながる）背景
% 2. 何が問題なのか？
% 3. この論文で解決したいことは何か（目的）？
% 4. 掲げた目的が達成されると何が嬉しいのか（既存手法では何故ダメなのか）？
% 5. 提案手法・システムの直感的なアイデア
% 6. 提案手法・システムの具体的な説明
% （より詳細な説明は別の章で記すので，項目5を実現するための要素を書く程度でよい）
% 7. 研究で得られた知見（どんな実験をして，その結果，どんな知見が得られたか）
%
% 以下は例です．
% > 1. ウェブには正しい情報もあれば間違った情報もある．相当数の人が誤ったウェブ情報を鵜呑みにしている．
% > 2. どのウェブ情報が正しいか間違っているかについて，一般ユーザが常に注意を払うことは難しい
% > 3. 本稿では，信憑性が怪しい情報に対するユーザの注意力を高める方法を提案する
% > 4. 既存手法として，情報の信憑性を自動的に分析する方法が提案されているが，すべての情報を自動分析することは困難．最終的には人間に判断が委ねられるので，人間の力を高める必要がある．
% > 5. 積極的に証拠を調べる気にさせるメッセージをウェブ検索中に提示する
% > 6. 認知心理学のプライミング効果を利用して，批判的な情報の見方を促すメッセージを自動生成し，ウェブ検索時に表示する
% > 7. ユーザ実験の結果，通常のウェブ検索エンジンを用いるよりも，公的機関や学術機関が発行している情報を多く閲覧するようになった．
%
%
% #### ユーザ調査系の場合
% 1. （次に述べる問題につながる）背景
% 2. 何が問題なのか？
% 3. この論文で明らかにしたいことは何か（リサーチクエスチョンRQは何か）？
% 4. 掲げたRQが明らかになると何が嬉しいのか（既存研究の限界は何か）？
% 5. RQを明らかにするための調査・実験のアイデア
% 6. 提案する調査方法・実験方法の具体的な説明
% （より詳細な説明は別の章で記すので，項目5を実現するための要素を書く程度でよい）
% 7. 研究で得られた知見（調査・実験をした結果，どんな知見が得られたか）
%
% 以下は例です：
% > 1. ウェブには正しい情報もあれば間違った情報もある．
% > 2. 相当数の人が誤ったウェブ情報を鵜呑みにしている．
% > 3. どんなユーザが誤った情報を取得しやすいのかが分かっていない
% > 4. 情報を批判的に調べる能力についての調査は主に高等教育機関の学生を対象に行われているが，それ以外の対象に対する調査はない．またそれら調査は汎用的な認知能力（批判的思考能力）の調査が目的であり，ウェブ情報をどのように閲覧しているかの調査事例はほとんどない．
% > 5. アンケート調査と実際の検索行動分析を組み合わせた分析を行う．
% > 6. アンケート調査に回答した被験者に対して，ウェブ検索タスクをさせて行動を分析．
% > 7. 誤った情報を取得しやすいユーザは，積極的に情報を調べる意欲が乏しく，ウェブ検索をしても検索結果の上位しか閲覧しない．
%
% 各要素を書くときは，1段落で1主張となるようにすると読みやすくなります．
%
% 上記アドバイスを補足するとすれば，論文全体を通じて説明に用いる共通の例を導入することです．取り扱う問題やその重要性を分かりやすく伝えるために例を使うのは有効です．提案事項の詳細を述べる章でも同じ例を用いることで，読者は論文をスムーズに読み進めることができます．
%
%
% -----------------------------
% 補足（参考：how-to-xx/how-to-write-a-paper/1-学術論文の書き方.md 3.3節）
% -----------------------------
%
% ■ 行き詰まったら「5つの問い」に答えてみる
% 上の7項目と対応しますが，より短く覚えられる型として次のものがあります．書き出せなくなったら，この5問に一つずつ答えるつもりで書いてみてください．
%
% 1. What is the problem?（問題は何か）
% 2. Why is it interesting and important?（なぜ興味深く重要か）
% 3. Why is it hard?（なぜ難しいか．素朴なアプローチはなぜ失敗するか）
% 4. Why hasn't it been solved before?（なぜ未解決だったか．既存アプローチに何が足りないか）
% 5. What are the key components of my approach and results?（アプローチと結果の鍵は何か．限界も明記する）
%
% ■ 章の最後に「本論文の貢献」を箇条書きで置く
% イントロダクションの最後に，「本論文の貢献は次の3点である」と貢献を箇条書きで明示してください．その際，次の2点を守ります．
%
% * 各貢献は「検証可能な主張」として書く．読者が「本当か？」と確かめられる粒度にします．
%   （悪い例：「〜を改善した」／良い例：「〜を用いることで検索回数が有意に増加することを示した」）
% * 各貢献に，それを述べている章・節の番号を添える（例：「(2) 〜する手法を提案する（3章）」）
%
% 貢献に章番号を添えておけば，イントロダクション自体が論文全体の見取り図を兼ねます．その結果，
% 「2章では背景を述べ，3章では手法を述べ…」という中身のない章構成の羅列を書く必要がなくなります．
%
% ■ 査読者は，イントロダクションを読み終えた時点で採否の初期判断をほぼ下しています
% 残りの部分は，その判断を裏付ける証拠を探しながら読まれます．冒頭で「そこで試合終了」と書いたのは，この意味です．
%
% ■ 共通の例（running example）は，この章の末尾で導入します
% 導入位置は，イントロダクション末尾の小節とするか，独立した章（2章または3章）を立てるのが一般的です．
% 一度決めた例は使い捨てにせず，提案手法の章でも実験の章でも同じ例を育てていってください．

% ----------------------
% はじめに/introduction
% ----------------------
\section{はじめに}


% 2章：関連研究
% -----------------------------
% 執筆にあたっての留意点
% -----------------------------
%
% 関連研究の章では，この論文で提案・主張することの意義を「相対的」に位置づけるために，関連する研究を紹介します．決して「類似する研究」を並べるわけではありません．自分の論文をうまく位置づけるためにも，3つ程度の観点（節）を設けて，関連研究を紹介するのがよいでしょう．

% 例えば，「情報の信憑性を積極的にチェックするための動機付けを行う方法」に関する論文を書いているならば，関連研究を整理する観点として

% ・信憑性が怪しい情報を計算機がチェックする方法に関する研究
% ・人間が自分の力で信憑性をチェックできるようになることを支援する研究
% ・そもそも人間がどのように信憑性をチェックしているのかを明らかにする研究

% などが観点として考えられます．

% 各観点（節）では最低でも3報は関連する論文を引用し，それぞれの論文で何を提案しているのか，何が明らかにされたのかを1，2行程度でまとめてください．その上で，各節の最後で，それら論文と自分の論文の方向性の違いを主張してください．
%
%
% -----------------------------
% 補足（参考：how-to-xx/how-to-write-a-paper/1-学術論文の書き方.md 4章）
% -----------------------------
%
% ■ 各段落の主役は，先行研究ではなく「自分の研究との関係」です
% 「Aは〜した．Bは〜した．Cは〜した」と並べただけでは，この章の役割を果たしていません．
% 先行研究という地図の中に自分の研究を位置づけ，「既存研究では埋まっていない穴があり，本研究がそれを埋める」
% ことを示すのが目的です．各節では，そのグループの到達点と限界を述べた上で，節の最後に
% 「これらに対し本研究は〜の点で異なる」と差分を明示してください．
%
% ■ 先行研究を貶めてはいけません
% 先行研究への敬意は，あなたの論文の価値を下げません．むしろ比較対象が優れているほど，
% それを超えるあなたの貢献が際立ちます．また現実問題として，査読者がその先行研究の著者本人である
% 可能性は決して低くありません．「〜には限界がある」と書くのは構いませんが，
% 「〜は不十分である」「〜は誤っている」といった断定は，根拠を示せない限り避けてください．
%
% ■ 引用は，必ず自分で原典を読んでから行ってください
% 他の論文の関連研究の記述をそのまま孫引きすると，誤った要約をそのまま引き継ぐことになります．
% 少なくともアブストラクトと図，結論には目を通してから引用しましょう．
%
% ■ この章をどこに置くか
% 本テンプレートでは2章（イントロダクションの直後）に置いていますが，末尾（考察と結論の間）に置く流儀もあります．
% 判断基準は次の通りです．
%
% * 冒頭に置く：短く，かつ十分詳細に書ける場合．または，先行研究に対して早い段階で自分の立場を明確にしたい場合
% * 末尾に置く：イントロダクションで手短に触れておける場合．または，先行研究と十分に比較するために提案手法の技術的な説明が必要な場合
%
% 迷ったら，投稿先の学会の採録論文を数本見て，その分野の慣習に合わせるのが安全です．


% ----------------------
% 関連研究/Related Work
% ----------------------

\section{関連研究}


% 3章：提案手法
% -----------------------------
% 執筆にあたっての留意点
% -----------------------------
%
% 「提案事項の詳細」について触れる章の名前は内容に依ります．単に「提案手法」「提案システム」とする場合もあれば，「XXを行う手法」と具体的な名前をつけることもあります．また，章を複数に分割する場合もあります．
%
% 提案手法のアイデアや挙動を説明するときは，必ず例を使いながら説明してください．そうすることで，読者が内容を理解しやすくなります．ここで使う例は「イントロダクション」で使った例と同じものを使うのが望ましいです．
%
%
% -----------------------------
% 補足（参考：how-to-xx/how-to-write-a-paper/1-学術論文の書き方.md 3.4節）
% -----------------------------
%
% ■ 例は使い捨てにせず，一つの例を論文全体で育ててください
% 上に書いた「イントロダクションと同じ例を使う」ことの狙いです．場面ごとに新しい例を出すと，
% 読者は新しい概念が出てくるたびに新しい状況を理解し直さなければなりません．
% 一つの例を，定義 → 手法 → 実行例 → 実験と一貫して使えば，読者は認知負荷を抑えたまま読み進められます．
%
% ■ 論文全体のおよそ1/4を読み終えるまでに，新規性のある技術的貢献を提示してください
% 8ページの論文なら2ページ目の終わりごろ，12ページの論文なら3ページ目あたりが目安です．
% そこまで読んでも「何が新しいのか」が見えない論文からは，読者も査読者も離れていきます．
%
% ■ 前提知識は切り分けて，先に簡潔にまとめます
% 自分の貢献ではない記法・用語・既存技術の説明は，「準備」などの節にまとめます．
% これは，「どこからが自分の貢献なのか」の境界を読者に示す役割を持ちます．
%
% ■ トップダウンに書きます
% 読者が「この先どこへ向かうのか」を常に見通せるよう，全体像 → 詳細の順で書いてください．
% 細部を読み飛ばしても手法の概要が掴めるのが理想です．
%
% ■ 設計上の選択には，必ず「なぜそうしたのか」の理由を添えます
% 「Aを採用した」ではなく，「Bでは〜という問題が生じるため，Aを採用した」と書いてください．
% 理由の書かれていない設計は，読者からは「なんとなくそうした」ようにしか見えません．
%
% ■ 語るべきは「完成した手法の物語」であって，「そこに至った試行錯誤の経緯」ではありません
% 自分が苦労した順に書きたくなりますが，読者が知りたいのは完成した手法の筋道です．
% 本筋を中断させる細部（証明の詳細，実装上の細かい工夫など）は，付録への退避を検討してください．
%
% ■ 読者が再現できる詳しさと，本質を見失わない抽象度のバランスを取ります
% パラメータの値やデータ構造の詳細は必要ですが，それらに埋もれて「何をしている手法なのか」が
% 見えなくなってはいけません．まず手法の全体像を1段落で述べてから，細部に入ってください．


% ----------------------
% 提案手法
% ----------------------
\section{提案手法}

% 4章：実験
% -----------------------------
% 執筆にあたっての留意点
% -----------------------------
%
% 実験・調査の章は，提案事項や仮説を立証するための評価実験，調査の方法を述べる章です．名前のつけ方は内容や分野によって異なります．仮説の立証方法が実験であれば「実験」「ユーザ実験」など，アンケートなどの調査であれば「調査」などが章の名前として考えられます．理学などの分野では，単に「方法」とする場合もあります．
%
% 「実験」の章では，前章までに設定した仮説を検証するための流れを説明してください．一般的には以下の項目に関して，節を立てて説明するとよいでしょう：
%
% * 実験協力者（実験協力者の属性，集め方）
% * 実験手順（時系列に沿って説明）
% * 実験装置・環境（どのような装置を用いたのか，どのような環境で実験をおこなったか，比較手法は何か）
% * 評価・測定指標（仮説と対応づける）
%
% 重要なことは，論文の読者が同じ実験を再現できるよう，必要な情報を正確かつ十分に書くことです．再現できない実験・調査を書いている場合は，反証可能性がないので「トンデモ研究」と見なされます．
%
%
% -----------------------------
% 補足（参考：how-to-xx/how-to-write-a-paper/1-学術論文の書き方.md 3.5節，
%              how-to-xx/how-to-write-a-paper/3-国際会議論文の傾向と対策.md 1章）
% -----------------------------
%
% ■ 検証したいことを，RQ（リサーチクエスチョン）や仮説として明示してから実験を設計します
% 「RQ1: 〜か？」「仮説1: 〜だろう」と番号を振っておくと，結果・考察の章と対応づけやすくなります．
% 測定指標は，必ずどれかのRQ・仮説に紐づけてください．どのRQにも紐づかない指標は，測る必要がありません．
%
% ■ 実験設計では「何を測るか」と「何を示すか」を分けて考えます
% * 何を測るか：実行時間，重要なパラメータへの感度，スケーラビリティ（データ量や問題の複雑さ），主観評価 など
% * 何を示すか：絶対性能（実用に足るか），素朴な手法との比較，既存手法との比較，
%               提案手法の構成要素ごとの寄与（アブレーション）
%
% ■ 自分の手法を良く見せるためだけの実験をしてはいけません
% 有利な条件だけを選んだ実験や，弱すぎるベースラインとの比較は，査読者にすぐ見抜かれます．
% 比較対象は「現時点で最も強い相手」を選んでください．また，なぜそのベースラインを選んだのかの
% 妥当性は，関連研究の章で説明できるようにしておきます．
%
% ■ 人を対象とする実験・調査では，倫理的な手続きを必ず記述します
% 倫理審査（所属機関の研究倫理委員会など）の承認，実験協力者への説明と同意の取り方，謝礼の有無と金額，
% 取得したデータの匿名化と管理方法を，この章に書いてください．国際会議では必須の記述です．
%
% ■ 分析方法も「再現できる粒度」で書きます
% * 定量分析：用いる検定の名前と，その前提（正規性，等分散性など）を確認したかどうか
% * 定性分析：thematic analysisなどの手法名と手順．複数人でコーディングした場合はその一致度，
%             1人で行った場合はその妥当性の説明
%
% ■ 実験の実施前に，この章を書き上げてしまうことを勧めます
% 実験の手順・条件・測定指標を文章にすると，設計の穴（統制していない要因，仮説と対応しない指標など）が
% 実施前に見つかります．実験をやり直す手間を考えれば，先に書いておく方がはるかに安上がりです．

% ----------------------
% 実験・調査
% ----------------------
\section{実験}

% 5章：結果
% -----------------------------
% 執筆にあたっての留意点
% -----------------------------
% 結果の章は実験・調査で得られたデータや分析結果について述べる章です．結果の章はデータや分析結果のみを淡々と記し，その解釈や結果から得られた知見は「考察」の章で述べるようにしてください．論文は「事実と意見を分けること」が重要となります．結果の章は実験・調査から得られた事実のみを述べるパートとし，意見は考察の章で述べるという棲み分けを行います．
%
% 結果は，測定項目・データ（=図や表）ごとに段落を分けて記述します．一般的には，各段落は，以下の3つの要素を順に並べて構成します．

% 1. 何の目的のために，何を調べたのかの説明（例：実験協力者がどの程度タスクに集中できたかを調べるため，タスク中のページ閲覧時間を調査した）
% 2. 得られた図・表の解説（例：図Xは，実験条件とページ閲覧時間の関係を示している）
% 3a. 図や表から読み取れる客観的事実（例：図Xによると，条件Aよりも条件Bのほうがページ閲覧時間が大きいことがうかがえる）
% 3b. 統計的な裏付け（例：分散分析を行った結果，条件Bと条件Aの間に統計的有意差があることが確認された）

% 前章で提示した仮説を漏れなく立証できるよう，必要となる実験・調査結果を記述してください．
%
%
% -----------------------------
% 補足（参考：how-to-xx/how-to-write-a-paper/1-学術論文の書き方.md 3.5節・5章，
%              how-to-xx/how-to-write-a-paper/2-文書の書き方.md 1章）
% -----------------------------
%
% ■ 都合の悪い結果も隠さずに書きます
% 仮説が支持されなかった結果や，提案手法が既存手法に負けた結果も，事実として報告してください．
% 隠して後から見つかるより，正直に報告した上で考察の章で誠実に議論する方が，査読では確実に好印象です．
%
% ■ 統計的な裏付けは，p値だけでは不十分です
% 有意差の有無に加えて，効果量と信頼区間も報告してください．「有意差があった」だけでは，
% その差が実質的に意味のある大きさなのかが読者に伝わりません．
% 検定を行う際は，その前提（正規性など）を満たしているかも確認しておきます．
%
% ■ 各段落の第1文で，その段落が何の結果の話なのかを述べます
% 上の構成の「1. 何の目的のために，何を調べたのか」がこれに当たります．
% 段落の先頭に「〜を調べるために，〜を測定した」と置けば，読者は図表を見る前に段落の主題を把握できます．
%
% ■ 図表の体裁
% * 図表はページ（2段組ならカラム）の上部に配置します
% * 本文で最初に参照するページか，その次のページに置きます
% * 図中の文字サイズは，本文とほぼ同じ大きさにします（縮小して貼り付けた図は，文字が潰れがちです）
% * すべての図表を本文中から参照します．本文から参照されない図表を置いてはいけません
% * 図表は，キャプションと軸ラベルだけを見て内容が理解できる状態を目指します

% ----------------------
% 結果
% ----------------------


\section{結果}

% 6章：考察
% -----------------------------
% 執筆にあたっての留意点
% -----------------------------
% 結果に意味を与える章であり，イントロダクションに次いで重要な章です．考察の章の役割は，結果の章で示した実験結果，調査結果を踏まえて仮説の検証を行い，結論にいたるまでの議論を展開することにあります．結果の章では実験や調査で得られた結果を淡々と報告すると述べましたが，考察の章に入ってようやく結果の解釈を行い，著者の「意見」を述べることができます．
%
% 考察の章では以下の項目について記します．
%
% * 研究目的の振り返り
% * 結果の解釈と仮説の立証結果
% * 提案手法，研究手法の制約条件，限界
% * 提案手法がうまく機能した理由，うまく機能しなかった理由の考察
% * 結果から導かれる提言
%
% 考察は，論文の中で最も執筆に神経を使う章です．特に文章の論理展開には気をつけてください．実験・調査結果を踏まえて何らかの主張を行うときは，論理的な飛躍がないように注意しましょう．何らかの提言を行う場合は，立証された仮説や分析結果に関連している必要があります．実験・調査結果をよりどころとしない主張はしてはいけません．
%
%
% -----------------------------
% 補足（参考：how-to-xx/how-to-write-a-paper/1-学術論文の書き方.md 3.5節，
%              how-to-xx/how-to-write-a-paper/3-国際会議論文の傾向と対策.md 2章）
% -----------------------------
%
% ■ 結果の章の要約を繰り返すだけの章にしてはいけません
% 「図Xでは〜だった」と再掲するのではなく，「なぜそうなったのか」「それは何を意味するのか」を書きます．
% 結果の再掲は，議論を成立させるために必要な最小限にとどめてください．
%
% ■ 主張の一つ一つに，対応する結果があるかを確認します
% 考察で述べる主張は，結果の章のどの数値・どの図表に支えられているかを言えなければなりません．
% 「〜だと思われる」「〜であろう」と書きたくなったら，その根拠が結果の章にあるかを確認してください．
% なければ，それは考察ではなく単なる感想です．
%
% ■ 限界（Limitations）は，隠すのではなく自分から書きます
% 実験協力者の偏り，タスクの人工性，一般化できる範囲などを，自分から具体的に述べてください．
% 査読者が気づく限界を著者が書いていない場合，「気づいていない」か「隠している」と受け取られます．
% 自分の研究の限界を正確に把握できていることは，研究者としての能力の証明になります．
% なお，「限界を書く」ことと「自信のなさを表明する」ことは違います．
% 「〜という限界はあるが，本研究の主張は〜の範囲では成立する」と，主張の有効範囲を示す書き方をしてください．
%
% ■ HCI系の会議に投稿する場合
% 得られた知見を「今後の設計に対する示唆（design implications）」の形でまとめるのが定番の型です．
% 「本研究の結果は，〜のようなインタフェースを設計する際に〜すべきことを示唆している」といった形で書きます．

% ----------------------
% 考察
% ----------------------
\section{考察}


% 7章：おわりに
% -----------------------------
% 執筆にあたっての留意点
% -----------------------------
%
% おわりに（結論）の章は，論文を読み終えた読者の頭に何を残すかを決める章です．
% 分量は多くありませんが，アブストラクトと並んで「そこだけ読む読者」が最も多い章でもあります．
% 参考：how-to-xx/how-to-write-a-paper/1-学術論文の書き方.md 3.6節
%
% 以下の順に書くとよいでしょう．
%
% 1. 本研究で何をしたか（1〜2文）
% 2. 何がわかったか（定量的な結果を使って具体的に）
% 3. その知見が持つ意味
% 4. 今後の課題（Future Work）
%
% ■ アブストラクトやイントロダクションのコピーにしてはいけません
% 最もやりがちな失敗です．結論の章は，イントロダクションで述べた主張を「実験を終えた今の言葉」で
% 述べ直す場です．たとえばイントロダクションで「検索行動を大きく変える」と書いたのであれば，
% 結論では「セッションあたりの検索回数が◯倍に増加することを確認した」と，
% 得られた数値を使って具体的に言い直します．
%
% ■ ここで新しい主張を初めて出してはいけません
% 結論に書いてよいのは，本文で根拠を示した内容だけです．
% 「本研究の結果は〜にも応用できるだろう」と書きたくなったら，それが結果に支えられているかを確認し，
% 支えられていないなら今後の課題として書いてください．
%
% ■ 今後の課題は，研究の価値の一部です
% 「〜が今後の課題である」と列挙して終わるのではなく，本研究の限界のうちどれが次に解くべき問題なのかを，
% 理由とともに述べてください．すでに取り組み始めている研究があれば，
% 「現在〜に取り組んでおり，予備的な結果は有望である」と書いておくと，
% そのテーマが自分のものであることを示せます．
%
% ■ 「今後の課題」が言い訳の場にならないようにしてください
% 「時間がなかったため〜は実施できなかった」といった書き方は，読者にとって何の情報にもなりません．
% 「〜を明らかにするには〜の実験が必要であり，これを次の課題とする」と，次の研究の設計として書きます．

% ----------------------
% おわりに/conclusion
% ----------------------
\section{おわりに}




% -------------------------
% 参考文献リスト
% -------------------------
\markboth{参考文献}{参考文献}
\bibliography{reference}
\bibliographystyle{ipsjsort}


\end{document}
